\bookchapter{Capítulo 32: Manual de primeros auxilios para la seguridad del móvil}

\chapterintro{Este capítulo es un «plan de emergencia». Ojalá nunca lo necesites, pero si de verdad pierdes el móvil, te lo roban o sufres una estafa, este manual te ayudará a reducir al máximo las pérdidas; y cuando llegan desastres como terremotos o lluvias torrenciales, el móvil también es una herramienta para recibir alertas, enviar tu ubicación, pedir socorro y mantener el contacto. Te recomendamos leerlo entero ahora mismo y guardar capturas de los pasos clave, o imprimirlo y dejarlo en casa: porque cuando ocurra puede que ya no tengas el móvil a mano, y en pleno desastre puede quedarse sin batería o sin señal.}

\booksection{Escenario cero: terremoto / tiempo extremo — cómo el móvil te ayuda a salvar la vida}

Aquí se protege tu integridad física, no tus cuentas ni tu dinero. La configuración de las alertas se explica en el Capítulo 13, «Alertas del tiempo, alertas de desastres naturales y alertas sísmicas»; aquí solo veremos qué hacer cuando suene la alarma.

\textbf{Suena la alerta de terremoto: primero refúgiate, luego evacúa, y no te quedes grabando vídeos}\par

\begin{enumerate}
  \item Refúgiate en el lugar más cercano: agáchate, cúbrete y sujétate (bajo una mesa, en la esquina de un muro de carga o junto a un mueble sólido), protegiendo cabeza y cuello; aléjate de los cristales, las fachadas exteriores y las lámparas colgantes.
  \item No salgas corriendo del edificio a ciegas: el temblor principal suele durar solo entre diez y unas decenas de segundos, y las escaleras pueden ser aún más peligrosas; los ascensores, jamás.
  \item Cuando el temblor pare, evacúa con orden: usa las escaleras, atento a las réplicas y a los objetos que puedan caer; dirígete a un espacio abierto, lejos de postes eléctricos, vallas publicitarias y edificios en mal estado.
  \item Usa el móvil solo para lo esencial: enviar un mensaje breve a tu contacto de emergencia («Estoy en XX, estoy bien/lesionado, necesito ayuda») gasta menos batería que una llamada larga; si puedes enviar SMS o WeChat, prioriza el texto.
  \item Sin señal: consulta el Capítulo 7, SOS por satélite / mensajes satelitales (si tu modelo lo admite); o sube a un punto alto o un espacio abierto e inténtalo de nuevo; en algunos casos, las conversaciones por Bluetooth sin internet también ayudan.
  \item Desactiva las funciones no esenciales: en una emergencia, activa el modo avión y enciéndelo por intervalos, apaga el escaneo Bluetooth y baja el brillo; así aguantarás unas horas más. Reserva la última barra de batería para dar señales de vida y para la información de rescate.
\end{enumerate}

\begin{tcbwarn}{Lluvias torrenciales / tifón / crecidas repentinas: actúa ya en la fase de alerta}
\begin{itemize}
  \item Si ves una alerta roja: reduce al máximo salir de casa y aléjate de espacios subterráneos, puentes bajo nivel, cauces de ríos y pendientes pronunciadas.
  \item Si quedas atrapado en la calle: sube a las plantas altas de un edificio sólido y no intentes cruzar el agua a la fuerza (el agua acumulada puede esconder alcantarillas abiertas y cables con corriente); al llamar al 119/110, indica con claridad los puntos de referencia y el nivel del agua.
  \item Antes de viajar: confirma el riesgo con la app del tiempo y las cuentas oficiales de alertas de tu zona; en viajes largos en coche, evita los tramos de montaña junto al agua o al borde de los precipicios.
  \item Acuerdos entre familiares: decid de antemano «si perdemos el contacto, vale la última ubicación enviada por WeChat o SMS», para evitar que cada uno llame frenéticamente y agote la batería de todos.
\end{itemize}

\end{tcbwarn}

\begin{tcbpractice}{4 cosas que conviene tener hechas de antemano}
\begin{enumerate}
  \item Activa las alertas sísmicas del sistema o una app de alertas oficial, y permítele saltarse el modo silencio.
  \item Fija 2 o 3 «contactos de emergencia» al principio de la agenda para que puedan aparecer en la pantalla de bloqueo (con la ficha médica de emergencia de iOS/Android o la información en pantalla de bloqueo).
  \item Ten una batería externa en casa y en la mochila; antes de un viaje largo, carga el móvil a tope y lleva el cable.
  \item Memoriza al menos un número que puedas marcar sin agenda (el móvil de un familiar) y escríbelo en un papel para la cartera.
\end{enumerate}

\end{tcbpractice}

\booksection{Escenario uno: has perdido el móvil o te lo han robado — los 30 minutos de oro}

Los primeros 30 minutos tras perder el móvil son el «tiempo de oro»; si actúas bien, salvarás tu dinero y tu privacidad. Sigue estos pasos en orden y no te saltes ninguno:

\textbf{Procedimiento estándar de los primeros 30 minutos tras la pérdida}\par

\begin{enumerate}
  \item Llama a tu propio número desde el móvil de otra persona (minuto 1): quizá solo se te ha caído cerca y, si la contesta alguien de buen corazón, podrás recuperarlo.
  \item Bloquea el móvil a distancia (minutos 2-5): desde el teléfono de otra persona o un ordenador, accede al servicio de búsqueda: iPhone: iCloud.com > Buscar > elige tu dispositivo > marcar como perdido
  \item Marcas Android: Huawei cloud.huawei.com / Xiaomi i.mi.com / OPPO cloud.oppo.com, entre otras
  \item Bloquea la tarjeta SIM (minutos 5-10): llama al teléfono de atención al cliente de tu operadora (China Mobile 10086 / China Unicom 10010 / China Telecom 10000), informa de la pérdida del móvil y pide suspender el servicio de la SIM. Este paso es el más crítico: si el estafador se queda con tu SIM, recibirá los códigos de verificación y podrá restablecer todas tus contraseñas.
  \item Congela las cuentas de pago (minutos 10-20): Alipay: llama al 95188 para reportar la pérdida, o inicia sesión en Alipay desde otro móvil > Yo > Ajustes > Centro de seguridad > Congelación de emergencia
  \item WeChat: tras denunciar al 110, contacta con el servicio de atención de WeChat, o entra en el Centro de seguridad de WeChat > Congelar cuenta
  \item App del banco: llama al teléfono de atención al cliente que aparece en la parte trasera de tu tarjeta y pide congelar la banca online
  \item Avisa a familia y amigos (minutos 20-25): publica en tu perfil de WeChat o pide a tu familia que difunda el aviso de que has perdido el móvil, para evitar que los estafadores usen tu identidad contra otras personas.
  \item Denuncia (minutos 25-30): llama al 110 o acude a la comisaría más cercana a presentar la denuncia y consigue el justificante. Lo necesitarás después para el seguro y para posibles disputas con el banco.
\end{enumerate}

\booksection{Escenario dos: te han estafado — las 72 horas de oro para reaccionar}

Al descubrir que te han estafado, lo más importante es mantener la calma. Decidir en medio del pánico provoca segundas pérdidas. Sigue estos pasos:

\begin{enumerate}
  \item Detén de inmediato cualquier transferencia: diga lo que diga el estafador, no envíes ni un céntimo más. «Transfiere otro depósito de garantía y te devolvemos lo anterior» es el clásico discurso de la doble estafa.
  \item Conserva todas las pruebas: ¡no borres los chats! ¡No borres el registro de llamadas! Haz capturas de los chats, las transferencias y los datos de la cuenta del estafador. Son las pruebas clave para la denuncia y para intentar recuperar el dinero.
  \item Llama de inmediato al 96110: es la línea nacional contra las estafas. El operador te guiará para pedir una parada de pago urgente: si la transferencia es reciente (normalmente dentro de las 24 horas), el banco puede congelar la cuenta del estafador e interceptar los fondos.
  \item Llama al 110 para denunciar: lleva el móvil, tu documento de identidad y los comprobantes de las transferencias a la comisaría más cercana para declarar.
  \item Contacta con el banco o la plataforma de pago: explica la estafa y pide congelar las cuentas implicadas. Si te han hecho cargos fraudulentos en la tarjeta de crédito, pide al banco que active el proceso de impugnación del pago.
  \item Cambia todas las contraseñas: desde otro dispositivo seguro, cambia la contraseña de tu Apple ID o cuenta de marca, WeChat, Alipay y las apps bancarias, y activa la verificación en dos pasos.
\end{enumerate}

\booksection{Escenario tres: el móvil se ha infectado con un virus o un troyano — ¿qué hacer?}

Si tu móvil presenta estos síntomas —anuncios que aparecen sin motivo, apps extrañas que se instalan solas, la factura del teléfono se dispara o el consumo de datos es anormal—, lo más probable es que tenga malware:

\begin{itemize}
  \item Primer paso: entra en modo seguro — en un Android, mantén pulsado el botón de encendido para apagar; luego mantenlo pulsado para encender y, en cuanto vibre, mantén pulsado el botón de bajar volumen sin soltarlo hasta entrar en modo seguro (abajo a la izquierda verás el texto «Modo seguro»). En este modo todas las apps de terceros se desactivan, lo que facilita la investigación.
  \item Segundo paso: desinstala las apps sospechosas — en modo seguro, entra en [Ajustes] > [Gestión de aplicaciones], ordena por fecha, localiza las apps instaladas recientemente que no reconozcas y desinstálalas.
  \item Tercer paso: instala una app de seguridad y analiza — descarga «Tencent Mobile Manager» o «360 Mobile Guard» y haz un análisis completo del sistema.
  \item Cuarto paso: restaurar de fábrica (último recurso) — si nada de lo anterior funciona, haz una copia de seguridad de los datos importantes y restaura el móvil de fábrica. Ojo: el restablecimiento borra todos los datos; antes copia sin falta las fotos, los contactos y otros archivos importantes.
  \item Un iPhone prácticamente no se infecta: el sistema de aislamiento (sandbox) de iOS dificulta que el malware sobreviva. Si tu iPhone se comporta de forma extraña, lo más probable es que sea un error de alguna app o de la configuración, no un virus. Restaurar de fábrica resuelve el 99 \% de los problemas.
\end{itemize}

\booksection{Escenario cuatro: te han robado una cuenta — recupera el control rápido}

\begin{enumerate}
  \item Cambia la contraseña de inmediato: si todavía puedes entrar en la cuenta, cámbiala en cuanto descubras el robo. La nueva debe ser totalmente distinta de la antigua.
  \item Expulsa los demás dispositivos: en los ajustes de seguridad de la cuenta, busca «gestión de dispositivos conectados» y cierra la sesión de todos los que no reconozcas.
  \item Revisa el historial de accesos: mira los inicios de sesión de los últimos 7 días; si hay conexiones desde otros lugares, la cuenta está efectivamente comprometida.
  \item Activa la verificación en dos pasos: si aún no la tienes, actívala ahora mismo. Es lo más eficaz para evitar que te roben la cuenta otra vez.
  \item Revisa los datos vinculados: comprueba si han cambiado el número de teléfono o el correo asociados a la cuenta. Si es así, recupérala con la función de apelación de la cuenta.
  \item Revisa los ajustes de pago: si te han robado WeChat o Alipay, comprueba que no hayan manipulado los pagos sin contraseña ni los cobros automáticos, y que no hayan añadido tus tarjetas a cuentas desconocidas.
\end{enumerate}

\booksection{Escenario cinco: recibes mensajes de chantaje — no temas, no pagues}

Hay dos chantajes típicos: (1) «tengo vídeos o fotos íntimas tuyos; si no pagas, se los envío a todos los contactos de tu agenda»; (2) «tu móvil está controlado por un hacker; si no pagas, lo bloqueo para siempre». Si te ocurre esto:

\begin{itemize}
  \item El 99,9 \% es estafa: en realidad no tienen ningún dato privado tuyo. Envían el mismo mensaje de chantaje a decenas de miles de personas; siempre hay quien se asusta y paga.
  \item No respondas ni pagues: responder equivale a decirles «este número está activo»; te marcarán como «objetivo chantajeable» y el acoso posterior será mayor.
  \item Haz capturas como prueba: y después bloquea al remitente.
  \item Cambia las contraseñas: si te preocupa, cambia la de tus cuentas principales para quedarte tranquilo.
  \item Si de verdad existe riesgo de filtración de tu intimidad: denuncia. La policía tiene experiencia en este tipo de casos.
\end{itemize}

\booksection{Escenario seis: estafas con caras y voces falsas por IA — ver ya no es creer}

Desde 2024, la IA puede clonar rostros y voces hasta el punto de que un estafador aparece en una videollamada «convertido» en tu familiar, tu jefe o tu amigo. «Mamá, me ha pasado algo, transfíreme dinero ya»: la persona del vídeo se ve y suena idéntica a tu hijo, pero es falsa. Es uno de los engaños más peligrosos y difíciles de detectar en la actualidad.

\textbf{Primer paso: detectar los fallos — las cuatro grietas típicas de la falsificación por IA}\par

\begin{itemize}
  \item Frecuencia de parpadeo anormal: una persona real parpadea entre 15 y 20 veces por minuto; las caras falsificadas por IA suelen parpadear demasiado poco o con un ritmo rígido.
  \item Labios y audio desincronizados: si te fijas bien, entre el movimiento de la boca y el sonido que oyes hay un desfase de unas décimas de segundo.
  \item Bordes del rostro borrosos: el contorno de la cara, la zona de las orejas y el borde del pelo suelen mostrar «rebabas» o parpadeos extraños.
  \item Luces y sombras incoherentes: los brillos y las sombras del rostro no coinciden con la dirección de la luz del fondo, o al girar la cara los rasgos «giran» de forma incorrecta.
\end{itemize}

\begin{tcbinfo}{Segundo paso: tres trucos para verificar quién está al otro lado}
\begin{itemize}
  \item Pide movimientos concretos: pídele que gire la cabeza, se tape la cara o se cubra la boca con la mano y la abra y cierre. La cara falsificada por IA se delata más fácilmente con oclusiones y giros amplios: se deforma o no sigue el movimiento.
  \item Haz preguntas que solo él conoce: inventa una pregunta falsa, como «¿en qué clínica llevamos la última vez al perro?»; si es un estafador, no sabrá responder o se enredará en evasivas.
  \item Devuelve la llamada al número conocido: corta el vídeo y llama directamente al número guardado en tu agenda. El estafador no controla el número real: en cuanto marcas, se descubre.
\end{itemize}

\end{tcbinfo}

\begin{tcbwarn}{Tercer paso: cómo denunciar tras confirmar la estafa}
\begin{enumerate}
  \item Detén las transferencias al instante: por mucha «urgencia» que haya, haz una pausa; casi todo el que te apura con «transfiere ya y no cuelgues» es un estafador.
  \item Conserva todas las pruebas: graba la pantalla y guarda con capturas la videollamada, los chats, la cuenta del estafador y los comprobantes de las transferencias: son el material clave para el análisis técnico.
  \item Llama al 96110, la línea contra estafas: indica que se trata de una estafa con «cara o voz falsificadas por IA»; el operador activará la parada de pago urgente.
  \item Llama al 110 para denunciar: acude con todas las pruebas a la comisaría a declarar y señala expresamente que se usó tecnología de falsificación por IA, para que la policía pueda clasificar bien el caso.
  \item Avisa a familia y amigos: el estafador puede usar tu vídeo falsificado para engañar a otras personas; publica un aviso por adelantado en tu perfil o en el grupo familiar.
\end{enumerate}

\end{tcbwarn}

\booksection{Escenario siete: el ataque de SIM swap (secuestro de SIM) — te han «robado» el número}

Es un ataque que mucha gente no conoce y que causa daños enormes. El SIM swap (secuestro de la tarjeta SIM) consiste en que el atacante, falsificando tus datos de identidad, se hace pasar por ti en la tienda del operador y emite una SIM nueva a tu nombre. Cuando lo consigue, la señal de tu móvil desaparece de golpe, y el teléfono del atacante, con esa tarjeta nueva, recibe todos tus códigos de verificación por SMS: Alipay, el banco, WeChat… todas las puertas quedan abiertas.

\begin{tcbinfo}{¿Cómo darte cuenta? Las tres señales de alerta del SIM swap}
\begin{itemize}
  \item El móvil muestra de pronto «Sin servicio» o «Sin SIM», aunque tú no hayas tocado la tarjeta y reiniciar no sirva de nada: suele ser el momento del secuestro.
  \item Recibes un SMS del operador con «cambio de tarjeta completado» o «trámite realizado con éxito», pero tú no has hecho ningún trámite.
  \item De pronto no te llegan los códigos de verificación, y tampoco el resto de SMS: porque los mensajes están siendo desviados a la tarjeta del atacante.
\end{itemize}

\end{tcbinfo}

\begin{tcbwarn}{Medidas de protección — ponle «candado» también a la SIM}
\begin{itemize}
  \item Configura el PIN de la SIM: ruta [Ajustes] > [Teléfono] > [PIN de la tarjeta SIM]. Una vez activado, cada vez que el móvil se reinicie o cambies de tarjeta habrá que introducir el PIN para poder usarla. Aunque el atacante se apodere de la tarjeta, no podrá utilizarla.
  \item Activa la verificación adicional del operador: entra en la app de China Mobile/Unicom/Telecom, busca «centro de seguridad» o «verificación de trámites», y activa el aviso por SMS cuando alguien pida reemplazar tu tarjeta sin ser tú y la doble comprobación de contraseña para los trámites. Así, cuando alguien intente emitir tu tarjeta, el empleado exigirá una verificación de identidad adicional.
  \item No dependas solo del SMS en las cuentas importantes: para cuentas de alto valor como el banco o el Apple ID, usa preferentemente apps de autenticación (Google Authenticator, Microsoft Authenticator) o llaves de hardware, en lugar de códigos por SMS.
  \item Añade marca de agua a las copias del documento de identidad: en las copias necesarias para trámites, escribe «solo para el trámite de XX»; los atacantes suelen suplantarte con datos de identidad filtrados.
\end{itemize}

\end{tcbwarn}

\textbf{Qué hacer tras el secuestro — la velocidad lo es todo}\par

\begin{enumerate}
  \item Llama de inmediato al servicio del operador para reportar la pérdida: China Mobile 10086 / Unicom 10010 / Telecom 10000; pide la suspensión inmediata de la línea y la congelación del trámite de reemplazo de SIM para impedir que el atacante siga usando tu número.
  \item Entra en tus cuentas clave desde otros dispositivos: con un ordenador o el móvil de un familiar, cambia enseguida las contraseñas del banco, Alipay, WeChat y Apple ID, y expulsa los dispositivos desconocidos.
  \item Congela las cuentas de pago: reporta la pérdida en Alipay llamando al 95188, congela la cuenta desde el Centro de seguridad de WeChat y pide al banco que bloquee la banca online.
  \item Ve a la tienda del operador con tu documento de identidad: preséntate en persona en una tienda oficial para emitir una tarjeta nueva y recuperar el uso de tu número.
  \item Denuncia: el SIM swap constituye un delito contra los datos personales de los ciudadanos; presenta la denuncia con todas las pruebas (registro de la hora de la suspensión, pérdidas económicas, historial del reemplazo de tarjeta).
  \item Revisión a fondo: tras recuperar el número, repasa uno a uno los accesos recientes y los movimientos de dinero para confirmar que no queda ningún rastro de suplantación.
\end{enumerate}

\begin{tcbwarn}{Prevenir es mejor que curar — 5 cosas del día a día}
\begin{enumerate}
  \item Contraseña de bloqueo + verificación en dos pasos: activa las dos
  \item Confirma que «Buscar mi móvil» está activado
  \item Ponle PIN a la SIM ([Ajustes] > [Teléfono] > [PIN de la tarjeta SIM]): si pierdes el móvil, nadie podrá usar tu tarjeta
  \item No instales apps de origen desconocido ni pulses enlaces extraños
  \item Haz copias de seguridad periódicas de los datos importantes (fotos, contactos, chats)
\end{enumerate}

\end{tcbwarn}

\begin{tcbpractice}{Manos a la obra: crea tu propia «tarjeta de contactos de emergencia»}
\begin{enumerate}
  \item Coge un papel y anota estos números: la atención al cliente de tu operadora (China Mobile 10086 / Unicom 10010 / Telecom 10000), la del banco, la línea contra estafas 96110, el teléfono de emergencias 110 y los teléfonos de tu familia.
  \item Anota el correo con el que registraste tu Apple ID o tu cuenta de marca y un medio de contacto alternativo.
  \item Guarda ese papel en la cartera o pégalo en la nevera: no lo guardes en el móvil, porque cuando ocurra algo puede que no lo tengas a mano.
  \item Compruébalo ahora mismo: ¿está activado «Buscar mi móvil»? ¿Has configurado el PIN de la SIM?
\end{enumerate}

\end{tcbpractice}

\begin{tcbsummary}{Resumen del capítulo}
\begin{itemize}
  \item Los 30 minutos de oro tras perder el móvil: llamar desde otro teléfono bloqueo remoto bloquear la SIM congelar los pagos avisar a familia y amigos denunciar
  \item Las 72 horas tras la estafa: detener las transferencias conservar las pruebas llamar al 96110 denunciar contactar con el banco cambiar las contraseñas
  \item Móvil infectado: modo seguro desinstalar apps sospechosas análisis de seguridad restaurar de fábrica (último recurso)
  \item Cuenta robada: cambiar contraseña expulsar dispositivos revisar accesos activar la verificación en dos pasos revisar datos vinculados revisar pagos
  \item Los mensajes de chantaje son estafa en el 99,9 \% de los casos: no respondas ni pagues; haz capturas y bloquea
  \item Estafas con cara o voz falsas por IA: vigila el parpadeo, los labios, los bordes y las luces; pide que gire la cabeza o se tape la cara, haz preguntas íntimas y devuelve la llamada al número conocido
  \item Secuestro por SIM swap: un «Sin servicio» repentino es señal de alerta; ponle PIN a la SIM, activa la verificación adicional del operador y usa apps de autenticación en las cuentas importantes
  \item Cinco hábitos de prevención diarios: bloqueo + 2FA, buscar el móvil, PIN de la SIM, no instalar apps de origen desconocido y copias de seguridad periódicas
\end{itemize}

\end{tcbsummary}
